\documentclass[11pt]{article}
\usepackage{amsmath,amssymb}
\title{A joint theta-network magnetic-current bound with unbalanced times}\author{}\date{}
\begin{document}\maketitle
\paragraph{Main-gap interface.}
The finite boundary nonlinear ward cancellation calculation gives the complete finite-boundary mass cancellation
and its three-dimensional shell accounting. The local noise ward and longitudinal cancellation calculation gives the local
coefficient cancellation of longitudinal vector powers. The remaining
cubic-pair graph has a genuinely joint path current. This note proves an
estimate on that current for ONE of the original scalar mixed-boundary
grids in the mixed boundary magnetic loop bound calculation, including arbitrarily unbalanced proper times. It does
not identify different staggered component grids, and does not assign an
unproved derivative factor to the original cubic vertices.

\paragraph{1. Exact original scalar network.}
Let $p_t$ be the original real killed/reflected kernel on one edge grid
(one Neumann and three Dirichlet directions) or the original ghost grid
(four Dirichlet directions). All THREE paths below use this same grid.
Use its actual constant-flux connection of strength $b=2/R^2$. For
surviving paths $\omega_i:x\to y$ of durations $t_i$, write $J_i$ for
the oriented magnetic current, with the same gauge for all paths.
The third path is neutral in the magnetic product; assigning it a
geometric current as an auxiliary variable does not change its measure.
Define the unnormalized joint current
\[
 \mathcal A_\theta(t_1,t_2,t_3)=
 \sum_{x,y}\int dP^{t_1}_{xy}(\omega_1)dP^{t_2}_{xy}(\omega_2)
                  dP^{t_3}_{xy}(\omega_3)\,(J_1-J_2)^2 .           \tag{1}
\]
It is exactly minus the second background derivative at zero of
$\sum_{x,y}p_{t_1}^{\epsilon}(x,y)
 p_{t_2}^{\epsilon}(y,x)p_{t_3}(x,y)$.
There is no normalization by return probability and no bath-law change.

Let $\mathcal A_{\rm loop}(s)$ denote the complete single-loop quantity
of the mixed boundary magnetic loop bound calculation. Its proof and the original product heat kernel give
\[
 \sup_{x,y}p_t(x,y)\le C k(t),\quad
 \mathcal A_{\rm loop}(s)\le C|\Omega|b^2 s^2 k(s),\quad
 k(t)=(1+t)^{-2}e^{-ct/R^2}.                                      \tag{2}
\]
Constants can be decreased in the exponential when combining bounds.

\paragraph{2. Use the shortest path as the current reference.}
If path three is not the shortest, bound its kernel by its supremum and
concatenate path one with the reversal of path two. Summing their common
endpoint produces EXACTLY a loop of length $t_1+t_2$. Thus
\[
 \mathcal A_\theta\le (\sup p_{t_3})
                        \mathcal A_{\rm loop}(t_1+t_2).           \tag{3}
\]
If path three IS shortest, do not take its possibly large kernel supremum.
Instead insert its geometric current before estimating:
\[
 J_1-J_2=(J_1-J_3)+(J_3-J_2).
\]
The square is bounded by twice the sum of the two squares. In each term
bound the UNUSED path kernel by its supremum; the other two paths again
concatenate to an exact loop. Consequently
\[
 \mathcal A_\theta\le2(\sup p_{t_2})\mathcal A_{\rm loop}(t_1+t_3)
             +2(\sup p_{t_1})\mathcal A_{\rm loop}(t_2+t_3).       \tag{4}
\]
This reference-path identity is valid even though path three is neutral:
only the algebraic current comparison changes, not its probability weight.
It uses no independent bridge-area estimate or assumption of comparable
times. Reversal is legitimate for the original real reversible kernel.

Sort the times as $m\le M\le T$. In (3)--(4) every nontrivial pair consists
of the shortest time and either $M$ or $T$. Their sums are comparable to
the longer member. Equations(2)--(4) therefore prove
\[
 \boxed{\mathcal A_\theta(t_1,t_2,t_3)
 \le C|\Omega|b^2\,
 \frac{T^2}{(1+M)^2(1+T)^2}
             e^{-c(M+T)/R^2}.}                                   \tag{5}
\]
This is a uniform-in-$R$ bound on the actual scalar joint current. It
retains both return factors of the two longer paths in the unbalanced
region, where a naive neutral-kernel supremum would lose the estimate.

\paragraph{3. The precise remaining derivative weight.}
The bound already gives two useful, distinct integral estimates:
\[
 \int_{(0,\infty)^3}\mathcal A_\theta\,dt_1dt_2dt_3
       \le C|\Omega|b^2 R^2\log(2+R),                             \tag{6}
\]
\[
 \boxed{\int_{(0,\infty)^3}
       {\mathcal A_\theta\over1+\max(t_1,t_2,t_3)}\,dt_1dt_2dt_3
       \le C|\Omega|b^2[\log(2+R)]^2.}                            \tag{7}
\]
Indeed on each of six ordering regions integrate $m$ first, producing
$M$, and then use
$\int_0^T M(1+M)^{-2}dM\le C\log(2+T)$.
For (7) the remaining integral is bounded by
$C\int_0^\infty T^2(1+T)^{-3}\log(2+T)e^{-cT/R^2}dT$,
which is $O(\log^2(2+R))$. Without the additional denominator it is
$O(R^2\log(2+R))$, proving (6).
Since $|\Omega|b^2\le4$, (7) has no volume factor for the actual physical
constant-flux source. However its denominator is an EXPLICIT UNPAID
vertex-derivative weight. It must not be inserted into a Wilson graph
without deriving it from the original covariant curl/ghost vertices and
all differentiated endpoint partners. The literal cubic Wilson vertex has
one plaquette curl; two such derivatives suggest this weight, but this
note does not prove the required current-weighted derivative estimate.

\paragraph{4. Verification and exact remaining joins.}
A finite-state recursion keeps all three original path
measures on the $R=3$ mixed grid with reflected coordinate one. It computes
first and second path-current moments without sampling. For durations
$(2,5,1),(1,2,7),(8,1,3),(16,2,1),(4,4,4),(1,24,2)$ in uniformized steps,
the exact inequalities (3)--(4) pass; the largest observed ratio to their
right side is $0.965$. These diagnostics verify algebra, not the analytic
asymptotic bounds.

The actual full vector/ghost coefficient also couples DIFFERENT staggered
component grids at plaquette vertices. Those kernels are not silently
identified here; a neutral auxiliary path on another grid cannot be used
in (4) without a proved common realization or an equivalent mixed-kernel
estimate. The component boundary conditions, local endpoint connectors,
spin insertions and changing physical projections remain part of that
join. The complete boundary Ward of the finite boundary nonlinear ward cancellation calculation and longitudinal cancellation
of the local noise ward and longitudinal cancellation calculation must also remain in force when supplying the derivative weight.
Thus (5) is a genuine joint-current building block; (7) quantifies the
specific marginal weight still needed. It is not a bound on the full
$\mathcal C_R''$, on the original nonlinear remainder, or on S94.
\end{document}
